\pagestyle{empty}

\noindent \textbf{\Large Adathordozó használati útmutató}

\vskip 1cm

Az adathordozón a \textit{bsc-degree-project} nevű projekt forráskódja és dokumentációja található, amely három fő részből áll: egy Node.js/Express alapú backend, egy React alapú webes frontend, valamint egy Expo alapú mobilalkalmazás.

\vskip 0.5cm

\noindent \textbf{Az adathordozó tartalma:}

\begin{itemize}
    \item \texttt{dolgozat.pdf} -- a szakdolgozat végleges változata,
    \item \texttt{latex/} -- a dolgozat \LaTeX{} forráskódja,
    \item \texttt{backend/} -- a szerveroldali alkalmazás forráskódja (Node.js, Express, TypeScript, MongoDB),
    \item \texttt{frontend/} -- a webes felhasználói felület forráskódja (React, Vite, TypeScript, Tailwind CSS),
    \item \texttt{mobile/} -- a mobilalkalmazás forráskódja (React Native, Expo),
    \item \texttt{README.md} -- telepítési és használati útmutató.
\end{itemize}

\vskip 0.5cm

\noindent \textbf{Szükséges szoftverek telepítése:}

Az alábbi szoftvereket globálisan kell telepíteni a rendszerre a projekt futtatása előtt.

\vskip 0.3cm

\noindent \textit{1. Node.js (18-as vagy újabb verzió)}

\noindent Letölthető: \texttt{https://nodejs.org}

\noindent Telepítés ellenőrzése:
\begin{verbatim}
node -v
npm -v
\end{verbatim}

\noindent \textit{2. MongoDB Community Edition}

\noindent Letölthető: \texttt{https://www.mongodb.com/try/download/community}

\noindent Telepítés után az adatbázis-szerver indítása:
\begin{verbatim}
mongod --dbpath /data/db
\end{verbatim}

\noindent A MongoDB alapértelmezetten a \texttt{mongodb://localhost:27017} címen fut.

\vskip 0.3cm

\noindent \textit{3. Expo Go (opcionális, fizikai eszközön való teszteléshez)}

\noindent Telepíthető az iOS App Store-ból vagy a Google Play áruházból. A mobilalkalmazás indítása után a terminálban megjelenő QR-kódot kell beolvasni.

\vskip 0.5cm

\noindent \textbf{Telepítés és indítás:}

A függőségek telepítése a projekt gyökérkönyvtárából egyetlen paranccsal elvégezhető:

\begin{verbatim}
npm run setup
\end{verbatim}

Ha újratelepítés szükséges (például verziókonfliktus esetén):

\begin{verbatim}
npm run reset
npm run setup
\end{verbatim}

Ezt követően az egyes komponensek az alábbi sorrendben indítandók el, mindegyik külön terminálablakban:

\begin{enumerate}
    \item A MongoDB szerver elindítása: \texttt{mongod}
    \item A backend szerver: \texttt{cd backend \&\& npm run dev} \hfill (port: \texttt{5000})
    \item A webes frontend: \texttt{cd frontend \&\& npm run dev} \hfill (port: \texttt{5173})
    \item A mobilalkalmazás: \texttt{cd mobile \&\& npx expo start}
\end{enumerate}

A backend működéséhez szükséges egy \texttt{backend/.env} konfigurációs fájl, amely az alábbi változókat tartalmazza:

\begin{verbatim}
PORT=5000
MONGO_URI=mongodb://localhost:27017/<adatbázisnév>
JWT_SECRET=<titkos_kulcs>
\end{verbatim}